% =============================================================================
%  Reporte (ES) — Dosis de neutrones secundarios en fantomas de agua
%  Fuente LaTeX (article). Compilar con pdfLaTeX + BibTeX.
%  Las figuras se cargan desde la carpeta figures/ mediante el comando \smartfig.
% =============================================================================
\documentclass[11pt,a4paper]{article}

\usepackage[utf8]{inputenc}
\usepackage[T1]{fontenc}
\usepackage[spanish,es-tabla]{babel}   % 'es-tabla' => tablas se llaman "Tabla"
\usepackage{lmodern}
\usepackage{microtype}
\usepackage[margin=2.3cm]{geometry}

\usepackage{amsmath,amssymb}
\usepackage{graphicx}
\usepackage{booktabs}
\usepackage{caption}
\usepackage{subcaption}
\usepackage{siunitx}
\usepackage{multirow}
\usepackage{authblk}
\usepackage[hidelinks]{hyperref}
\usepackage{cite}

\sisetup{detect-weight=true, detect-family=true}

% --- Figura "inteligente": usa el PNG si existe; si no, muestra un placeholder ---
\newcommand{\smartfig}[3][\linewidth]{%
  \IfFileExists{figures/#2}%
    {\includegraphics[width=#1]{figures/#2}}%
    {\fbox{\parbox[c][0.22\textheight][c]{0.92\linewidth}{\centering\ttfamily
       [figura pendiente] \detokenize{figures/#2} \\[4pt]\normalfont\small #3}}}%
}

\title{\textbf{Dosis de neutrones secundarios en fantomas de agua con fotones
monoenergéticos mediante simulación Monte Carlo}}

\author[1]{Jorge Ramírez López}
\author[2]{Valeria Catalina Siordia Ortíz}
\author[3]{Dr.\ Teodoro Rivera Montalvo (asesor)}
\author[3]{Lic.\ Alejandra Inzunza Lugo (asesora)}
\author[3]{Lic.\ César Samuel Romero Núñez (asesor)}
\affil[1]{Licenciatura en Física, División de Ciencias Básicas, CUCEI, Universidad de Guadalajara \\ \texttt{jorge.ramirez9331@alumnos.udg.mx}}
\affil[2]{Ingeniería Biomédica, Centro Universitario de Tlajomulco, Universidad de Guadalajara}
\affil[3]{Laboratorio de Dosimetría, CICATA Unidad Legaria, Instituto Politécnico Nacional}
\date{}

\begin{document}
\maketitle

\begin{abstract}
\noindent
Comprender la distribución espacial de dosis es esencial para evaluar cómo la
energía depositada por el haz alcanza el volumen tumoral minimizando la exposición
de los tejidos sanos circundantes. En particular, la generación de neutrones
secundarios por procesos fotonucleares constituye una fuente adicional de radiación
que, aunque de baja magnitud, puede afectar la dosis fuera del campo e incrementar
el riesgo de efectos biológicos no deseados. En este trabajo se caracterizó la
distribución espacial de dosis en un fantoma de agua mediante simulaciones Monte
Carlo en \texttt{Geant4}, con haces de fotones monoenergéticos de 6, 10, 12, 15 y
18~MeV. Para cada energía se realizaron 100 corridas de \(10^{6}\) partículas, y los
resultados se procesaron en \texttt{MATLAB} para obtener curvas de porcentaje de
dosis en profundidad (PDD), métricas globales (\(D_{\max}\) y su profundidad), mapas
2D/3D por componente e histogramas energéticos. La dosis total se descompuso en
contribuciones primaria y secundaria (electrones, positrones, fotones y
\textbf{neutrones}), con énfasis en la componente neutrónica. Las curvas PDD
reprodujeron el comportamiento físico esperado. Se identificó una componente
neutrónica de baja magnitud, creciente con la energía y espacialmente más difundida,
coherente con la producción fotonuclear. Esto evidencia la importancia de
cuantificar la componente neutrónica por su potencial impacto en la inducción de
cáncer y la necesidad de considerarla en estudios dosimétricos y en la planificación
de tratamientos de radioterapia.
\end{abstract}

\noindent\textbf{Palabras clave:} Monte Carlo, Geant4, PDD, fotoneutrones, dosimetría.

% =============================================================================
\section{Introducción}
El cáncer, entendido como el crecimiento anómalo de células malignas, puede tratarse
mediante radioterapia, es decir, el uso de radiación ionizante para tratar tumores
localizados. Su aplicación exige una caracterización precisa de la dosis, tanto
dentro como fuera del campo de tratamiento~\cite{khan2014radiation}. El porcentaje de
dosis en profundidad (PDD) describe cómo varía la dosis absorbida con la profundidad
cuando un haz incide sobre un paciente o maniquí~\cite{khan2014radiation}. Para ello
se emplearon simulaciones Monte Carlo, consideradas estándar de
oro~\cite{fielding2023montecarlo}, con el \textit{toolkit}
Geant4~\cite{agostinelli2003geant4}, que modela con exactitud y trazabilidad las
interacciones de la radiación con la materia.

En tratamientos con fotones de alta energía (\(\gtrsim 10~\mathrm{MV}\)) pueden
ocurrir reacciones fotonucleares que generan un campo de neutrones secundarios
(térmicos y rápidos), contaminando el haz terapéutico con dosis no
deseada~\cite{banaee2021}. Su impacto es significativo fuera del campo, donde los
sistemas de planeación comerciales tienden a subestimar esta
componente~\cite{kry2017tg158}. Cuantificar la contribución neutrónica es relevante
por su Eficacia Biológica Relativa (RBE)~\cite{johns1983physics}: los neutrones
presentan efectividad máxima para el daño complejo en el ADN cerca de
\(1~\mathrm{MeV}\)~\cite{baiocco2016origin}, por lo que aun a dosis bajas su
toxicidad motiva caracterizar la dosis física depositada~\cite{kry2017tg158}.

A pesar de esta evidencia, la normativa mexicana presenta una brecha regulatoria: la
NOM-033-NUCL-2016~\cite{nom033_2016} exige detectores de neutrones pero no establece
procedimientos operativos ante mediciones significativas, y ni ésta ni la
NOM-002-SSA3-2017~\cite{nom002ssa32017} detallan cómo integrar la dosis neutrónica en
la planificación clínica. Este trabajo evalúa la distribución espacial de dosis de
haces de fotones monoenergéticos (6, 10, 12, 15 y 18~MeV) en un fantoma de agua
voxelizado, descomponiendo la dosis en componentes primaria y secundaria con énfasis
en la neutrónica. El procesamiento se realizó en MATLAB R2025a sobre infraestructura
de cómputo de alto desempeño, en el marco del XXX Verano del Programa
Delfín~\cite{programadelfin2025} en el CICATA-Legaria del IPN.

% =============================================================================
\section{Marco teórico}
La \textbf{dosis absorbida} \(D\) es la energía depositada por unidad de masa
(\(D=\mathrm{d}\epsilon/\mathrm{d}m\)); su unidad SI es el gray
(Gy)~\cite{attix2004dosimetry}. El \textbf{porcentaje de dosis en profundidad}
normaliza la dosis a una profundidad \(d\) respecto a una de referencia (típicamente
\(d_{\max}\))~\cite{khan2014radiation}:
\begin{equation}
  \mathrm{PDD}(d,\mathrm{SSD}) = 100 \times
  \frac{D(d,\mathrm{SSD})}{D(d_{\mathrm{ref}},\mathrm{SSD})}.
\end{equation}

A energías superiores a \(\sim 10~\mathrm{MeV}\) los fotones pueden generar neutrones
por \textbf{reacciones fotonucleares}, proceso dominado por la Resonancia Gigante
Dipolar; para el oxígeno-16 del agua la sección eficaz tiene un pico cercano a
23~MeV~\cite{geant4physref}. Los fotoneutrones se moderan por colisiones con
hidrógeno hasta volverse térmicos (\(\sim 0.025~\mathrm{eV}\))~\cite{lakshmanan1982}.
Dado que la dosis no predice por sí sola el efecto biológico, se emplea la RBE, que
compara una radiación de prueba con una de referencia para el mismo
efecto~\cite{reeve2018}. Los neutrones producen secundarias de alta Transferencia
Lineal de Energía (LET), causando daño complejo y una RBE significativamente mayor
que 1~\cite{reeve2018}.

% =============================================================================
\section{Metodología}

\subsection{Configuración de Geant4}
Las simulaciones se ejecutaron en \texttt{Geant4 v11.3.2} en modo multihilo. La
física electromagnética se describió con \texttt{G4EmStandardPhysics}; la producción
de neutrones por reacciones fotonucleares se habilitó con \texttt{G4EmExtraPhysics}
(proceso \texttt{GammaNuclear}) junto con \texttt{G4HadronPhysicsQGSP\_BIC\_HP} y
\texttt{G4HadronElasticPhysics} para la dispersión elástica e interacciones
inelásticas de neutrones con precisión en el rango térmico a intermedio (modelo
\textit{High Precision}). Es una versión optimizada de la lista
\texttt{QGSP\_BIC\_HP}~\cite{geant4qgspbic}, sin procesos de decaimiento ni iones por
su baja relevancia en haces de fotones de 6--18~MeV.

\subsection{Entorno de simulación}
El entorno se implementó con \texttt{G4VUserDetectorConstruction}, definiendo un
fantoma cúbico de agua (\texttt{G4\_WATER}) de 40~cm de lado. El volumen se
discretizó en una malla de \(27\times27\times200\) vóxeles; dado que 40~cm no es
divisible entre los 15~mm nominales, el paso real resultó de \(400/27=14.8148\)~mm en
\(X/Y\) y exactamente 2~mm en \(Z\) (volumen \(\approx 439~\mathrm{mm}^3\) por
vóxel).\footnote{El diseño nominal era \(15\times15\times2~\mathrm{mm}^3\); el
redondeo a 27 vóxeles fija el paso real citado.} Un esquema se muestra en la
Fig.~\ref{fig:esquema} y una vista del entorno en la Fig.~\ref{fig:entorno}.

\begin{figure}[t]
  \centering
  \smartfig[0.8\linewidth]{esquema_geometrico.png}{Esquema geométrico del entorno.}
  \caption{Esquema geométrico del entorno de simulación en Geant4.}
  \label{fig:esquema}
\end{figure}

\begin{figure}[t]
  \centering
  \smartfig[0.85\linewidth]{visualizador_18MeV.png}{Trazas del haz en el fantoma.}
  \caption{Trazas de un haz angosto de 10 fotones primarios de 18~MeV interactuando
  con el fantoma de agua. La malla mostrada es de baja resolución, sólo ilustrativa.}
  \label{fig:entorno}
\end{figure}

\subsection{Generación del haz y ejecución}
El haz se configuró con \texttt{G4VUserPrimaryGeneratorAction} y se controló mediante
macros: fotones monoenergéticos de 6, 10, 12, 15 y 18~MeV incidiendo
perpendicularmente desde una fuente puntual centrada (campo plano ideal). Las
acciones por corrida se gestionaron con \texttt{G4UserRunAction} y
\texttt{G4VUserActionInitialization}, registrando la energía depositada por el haz
primario y por las secundarias. Los datos se guardaron en \texttt{.csv} (uno por
ejecución) con todos los vóxeles de dosis no nula y el desglose por componente. Cada
simulación usó \(10^{6}\) partículas por ejecución, con 100 ejecuciones
independientes por energía. La ejecución en lotes se automatizó con \texttt{bash} y
archivos macro.

\subsection{Procesamiento y estimación de incertidumbre}
El procesamiento se automatizó en \texttt{MATLAB}~\cite{matlab2024}: lectura de los
CSV crudos, cálculo de la dosis media 3D por vóxel, su desviación estándar y la curva
PDD; métricas por corrida y por energía (\(D_{\max}\), su profundidad, medias,
desviaciones y \(\epsilon=100\,\sigma/\mu\)); histogramas; mapas 2D/3D por componente;
y exportación de archivos promediados y tablas de resumen.

Para cada vóxel con dosis media \(\bar D>0\) se estimó la incertidumbre estadística
de la media como \(\epsilon_{\text{vox}} = 100\,(\sigma_{\text{vox}}/\sqrt{N})/\bar
D\), reportando su promedio espacial \(\overline{\epsilon}_{\text{vox}}\) sobre los
vóxeles con \(\bar D>0\). El uso del SEM relativo es consistente con la práctica en
códigos Monte Carlo consolidados~\cite{mcnp_v4c2000}. Para métricas globales por
energía se reporta \(\epsilon=100\,\sigma/\mu\), que cuantifica la variabilidad entre
corridas. Ambos caracterizan aspectos complementarios de la incertidumbre.

% =============================================================================
\section{Resultados}

\begin{figure}[t]
  \centering
  \smartfig[0.85\linewidth]{PDD_comparacion.png}{Curvas PDD por energía.}
  \caption{Curvas de PDD para las energías consideradas.}
  \label{fig:pdd}
\end{figure}

\begin{table}[t]
  \centering
  \caption{Resumen de parámetros dosimétricos por energía. Medias \(\mu\) y errores
  relativos \(\epsilon=100\,\sigma/\mu\) estimados a partir de 100 corridas por
  energía.}
  \label{tab:resumen}
  \begin{tabular}{l S[table-format=2.2] S[table-format=1.2] S[table-format=2.2] S[table-format=2.2] S[table-format=1.2] S[table-format=1.2]}
    \toprule
    \multirow{2}{*}{\textbf{Energía}} &
      \multicolumn{2}{c}{\textbf{Dosis máxima}} &
      \multicolumn{2}{c}{\textbf{Profundidad de \(D_{\max}\)}} &
      \multicolumn{2}{c}{\textbf{Dosis total}} \\
    \cmidrule(lr){2-3}\cmidrule(lr){4-5}\cmidrule(l){6-7}
      & {\(\mu\) [\si{\nano\gray}]} & {\(\epsilon\) [\%]}
      & {\(\mu\) [\si{\milli\metre}]} & {\(\epsilon\) [\%]}
      & {\(\mu\) [\si{\micro\gray}]} & {\(\epsilon\) [\%]} \\
    \midrule
    6 MeV  & 11.03 & 0.59 & 26.26 & 10.10 & 1.61 & 0.08 \\
    10 MeV & 15.59 & 0.50 & 43.88 &  6.47 & 2.37 & 0.09 \\
    12 MeV & 17.79 & 0.42 & 52.58 &  6.12 & 2.74 & 0.11 \\
    15 MeV & 21.12 & 0.35 & 65.72 &  4.17 & 3.28 & 0.12 \\
    18 MeV & 24.39 & 0.39 & 77.26 &  3.94 & 3.82 & 0.11 \\
    \bottomrule
  \end{tabular}
\end{table}

\begin{table}[t]
  \centering
  \caption{Error relativo medio por vóxel de la dosis media
  (\(\overline{\epsilon}_{\text{vox}}=\mathrm{SEM}/\bar D\), en \%). Promedio espacial
  sobre vóxeles con dosis \(>0\).}
  \label{tab:err_voxel}
  \begin{tabular}{l S[table-format=1.2]}
    \toprule
    \textbf{Energía} & {\(\overline{\epsilon}_{\text{vox}}\) [\%]} \\
    \midrule
    6 MeV  & 5.53 \\
    10 MeV & 5.54 \\
    12 MeV & 5.49 \\
    15 MeV & 5.44 \\
    18 MeV & 5.36 \\
    \bottomrule
  \end{tabular}
\end{table}

\begin{figure}[t]
  \centering
  \begin{subfigure}[t]{0.32\textwidth}\centering
    \smartfig{Hist_DosisMax_6MeV.png}{6 MeV}\caption{6 MeV}\end{subfigure}\hfill
  \begin{subfigure}[t]{0.32\textwidth}\centering
    \smartfig{Hist_DosisMax_10MeV.png}{10 MeV}\caption{10 MeV}\end{subfigure}\hfill
  \begin{subfigure}[t]{0.32\textwidth}\centering
    \smartfig{Hist_DosisMax_12MeV.png}{12 MeV}\caption{12 MeV}\end{subfigure}
  \\[4pt]
  \begin{subfigure}[t]{0.32\textwidth}\centering
    \smartfig{Hist_DosisMax_15MeV.png}{15 MeV}\caption{15 MeV}\end{subfigure}\hfill
  \begin{subfigure}[t]{0.32\textwidth}\centering
    \smartfig{Hist_DosisMax_18MeV.png}{18 MeV}\caption{18 MeV}\end{subfigure}\hfill
  \begin{subfigure}[t]{0.32\textwidth}\end{subfigure}
  \caption{Histogramas de \textbf{dosis máxima} (nGy) por energía; valores de \(\mu\),
  \(\sigma\) y \(\epsilon\) en la Tabla~\ref{tab:resumen}.}
  \label{fig:hist-dosis}
\end{figure}

\begin{figure}[t]
  \centering
  \smartfig[0.85\linewidth]{Maps2D_15MeV.png}{Mapas 2D a 15 MeV.}
  \caption{Mapas de dosis 2D (proyección de máxima intensidad en XY) para el haz de
  15~MeV: dosis total, fotones primarios, electrones, positrones, fotones, neutrones
  y otras secundarias.}
  \label{fig:maps2d}
\end{figure}

\begin{figure}[t]
  \centering
  \begin{subfigure}[b]{0.32\textwidth}\centering
    \smartfig{Map3D_15MeV_Total.png}{Total}\caption{Dosis total}\end{subfigure}\hfill
  \begin{subfigure}[b]{0.32\textwidth}\centering
    \smartfig{Map3D_15MeV_eminus.png}{e-}\caption{Electrones}\end{subfigure}\hfill
  \begin{subfigure}[b]{0.32\textwidth}\centering
    \smartfig{Map3D_15MeV_eplus.png}{e+}\caption{Positrones}\end{subfigure}
  \\[4pt]
  \begin{subfigure}[b]{0.32\textwidth}\centering
    \smartfig{Map3D_15MeV_neutron.png}{n}\caption{Neutrones}\end{subfigure}\hfill
  \begin{subfigure}[b]{0.32\textwidth}\centering
    \smartfig{Map3D_15MeV_other.png}{otras}\caption{Otras}\end{subfigure}\hfill
  \begin{subfigure}[b]{0.32\textwidth}\end{subfigure}
  \caption{Distribución de dosis 3D para el haz de 15~MeV. Los vóxeles se visualizan
  si su dosis normalizada supera 25\,\% (1\,\% para neutrones).}
  \label{fig:maps3d}
\end{figure}

% =============================================================================
\section{Conclusiones}
Las curvas de PDD (Fig.~\ref{fig:pdd}) reproducen el comportamiento físico esperado:
región de acumulación, aumento de la profundidad de \(D_{\max}\) con la energía y
caída suave en la cola, lo que respalda la consistencia del modelo. Las métricas
globales muestran gran estabilidad: \(D_{\max}\) presenta errores relativos muy bajos
(\(\epsilon\approx 0.35\text{--}0.59\%\)), mientras la profundidad de \(D_{\max}\)
tiene mayor dispersión que decrece con la energía (\(\epsilon\) de \(\sim 10\%\) a
\(\sim 4\%\)), en concordancia con los histogramas (Fig.~\ref{fig:hist-dosis}) y la
Tabla~\ref{tab:resumen}. A nivel espacial, el SEM relativo vóxel a vóxel se mantuvo en
\(5.3\text{--}5.5\%\) (Tabla~\ref{tab:err_voxel}).

Las visualizaciones 2D/3D (Figs.~\ref{fig:maps2d} y \ref{fig:maps3d}) muestran que
las componentes cargadas (\(e^-/e^+\)) concentran la dosis cerca del eje del haz,
mientras que los neutrones secundarios aparecen más dispersos y a niveles bajos, con
mayor presencia relativa a energías altas. Aunque su contribución dosimétrica es
pequeña, su potencial relevancia radiobiológica justifica caracterizarla por separado
y, en el futuro, traducirla mediante magnitudes como LET y RBE. El trabajo constituye
un marco reproducible para comparar energías y componentes de dosis.

\subsection*{Limitaciones y trabajo futuro}
(1) Haces monoenergéticos e idealizados; (2) maniquí homogéneo y sin cabezal de
acelerador; (3) resolución de vóxel fija; (4) no se evaluó LET/RBE. Extensiones
propuestas: incorporar espectros clínicos y geometría del cabezal; estudiar
escenarios fuera de campo; calcular LET y estimar RBE para secundarios
(especialmente neutrones); y aplicar reducción de varianza para disminuir
\(\overline{\epsilon}_{\text{vox}}\) manteniendo el costo de cómputo.

\subsection*{Agradecimientos}
Al Programa Delfín y a los compañeros delfines; al CICATA-Legaria y al equipo del
Laboratorio de Dosimetría; y en especial al Dr.\ Teodoro, a Alejandra y a César por
su guía, paciencia y acompañamiento.

\bibliographystyle{unsrt}
\bibliography{references}

\end{document}
